\subsection{\texorpdfstring{AUTOMORPHISMS OF \(\mathbf{T}^n\)}{AUTOMORPHISMS OF T TO THE N}}

Let H be a closed subgroup of \(R^{n}\) and let \(\varphi\) be the canonical homomorphism of \(R^{n}\) onto the quotient group \(R^{n}/H\) . If f is a continuous homomorphism of \(R^{n}/H\) into a topological group G, then \(\dot{f} = f \circ \varphi\) is a continuous homomorphism of \(R^{n}\) into G which is periodic and has a group of periods containing H; conversely every periodic continuous homomorphism of \(R^{n}\) into G, whose group of periods contains H, is of this form.

In the case where \(H = Z^{n}\) , the quotient group \(R^{n}/Z^{n} = T^{n}\) is compact, and therefore every continuous homomorphism f of \(T^{n}\) into a topological group G is a strict morphism of \(T^{n}\) into G, provided that G is Hausdorff (Chapter III, §2, no. 8, Remark 1); \(\dot{f} = f \circ \varphi\) is a strict morphism of \(R^{n}\) into G; moreover, \(f(\mathbf{T}^{n}) = \dot{f}(\mathbf{R}^{n})\) is a compact subgroup of G, isomorphic to a group \(T^{p}\) ( \(o \leqslant p \leqslant n\) ).

In particular, we see that the only continuous homomorphism of \(T^{n}\) into a group \(R^{m}\) is the zero mapping, since \(\{0\}\) is the only compact subgroup of \(R^{m}\) .

Let us apply this to continuous homomorphisms of \(\mathbf{T}^n\) into a group \(\mathbf{T}^p\) ; if \(f\) is such a homomorphism, \(\varphi\) the canonical homomorphism of \(\mathbf{R}^n\) onto \(\mathbf{T}^n\) , then \(f \circ \varphi\) is a continuous homomorphism of \(\mathbf{R}^n\) into \(\mathbf{T}^p\) ; hence (no. 3, Proposition 3) if \(\psi\) denotes the canonical homomorphism of \(\mathbf{R}^p\) onto \(\mathbf{T}^p\) , there exists a linear mapping \(u\) of \(\mathbf{R}^n\) into \(\mathbf{R}^p\) such that \(f \circ \varphi = \psi \circ u\) . If \(\boldsymbol{x} \in \mathbf{Z}^n\) , \(f(\varphi(\boldsymbol{x}))\) is the identity element of \(\mathbf{T}^p\) , so that we must have \(u(\boldsymbol{x}) \in \mathbf{Z}^p\) , i.e., \(u(\mathbf{Z}^n) \subset \mathbf{Z}^p\) . Conversely, if \(u\) is any linear mapping of \(\mathbf{R}^n\) into \(\mathbf{R}^p\) which satisfies this condition, then \(\psi \circ u\) is a periodic continuous homomorphism of \(R^n\) into \(T^p\) , whose group of periods contains \(Z^n\) , and therefore defines a continuous homomorphism of \(T^n\) into \(T^p\) .

Consider under what conditions \(f\) is an isomorphism of \(\mathbf{T}^n\) onto a subgroup of \(\mathbf{T}^p\) . First of all, \(u\) must be an injective mapping of \(\mathbf{R}^n\) into \(\mathbf{R}^p\) , otherwise the vector subspace \(\overline{u}^{-1}(\mathrm{o})\) would contain points \(x \neq 0\) arbitrarily close to the origin, and at such a point we should have

\[
f (\varphi (x)) = f (\varphi (0)) \quad \text { and } \quad \varphi (x) \neq \varphi (0),
\]

contrary to hypothesis. This condition implies that \(p \geqslant n\) . The image \(u(\mathbf{Z}^n)\) is then a discrete subgroup of rank \(n\) of the group \(\mathbf{Z}^p\) ; the invariant factors of \(u(\mathbf{Z}^n)\) with respect to \(\mathbf{Z}^p\) (§ 1, no. 1) must all be equal to 1, otherwise there would exist a point \(x \in \mathbf{Z}^n\) and an integer \(k > 1\) such that \(u(k^{-1}\mathbf{x}) \in \mathbf{Z}^p\) and \(k^{-1}\mathbf{x} \notin \mathbf{Z}^n\) , hence \(f(\varphi(k^{-1}\mathbf{x})) = f(\varphi(0))\) and \(\varphi(k^{-1}\mathbf{x}) \neq \varphi(0)\) , contrary to hypothesis. Conversely, if this condition is satisfied, \(u(\mathbf{R}^n) \cap \mathbf{Z}^n = u(\mathbf{Z}^n)\) , and \(f\) is an isomorphism of \(\mathbf{T}^n\) onto \(u(\mathbf{R}^n)/u(\mathbf{Z}^n)\) .

If we apply this argument to the case p = n, we have the following proposition:

PROPOSITION 5. Every isomorphism of the topological group \(\mathbf{T}^n\) onto one of its subgroups is an automorphism of \(\mathbf{T}^n\) which is obtained by passing to the quotient from a linear mapping \(u\) of \(\mathbf{R}^n\) onto itself which, restricted to \(\mathbf{Z}^n\) , is an automorphism of \(\mathbf{Z}^n\) .

Equivalently (§ 1, no. 1), if \(u(\boldsymbol{e}_i) = \sum_{j=1}^{n} a_{ij} \boldsymbol{e}_j\) , the \(a_{ij}\) must be integers such that \(\det(a_{ij}) = \pm 1\) . In particular, for \(n = 1\) :

PROPOSITION 6. The only isomorphisms of the topological group T onto one of its subgroups are the identity mapping and the symmetry \(x \to -x\) .

